\documentclass{article}
\usepackage[utf8]{inputenc}
\usepackage[russian]{babel}
\usepackage{amsmath}

\title{Силы трения}

\begin{document}
\maketitle

2D случай. Есть N тел и M контактов.

Для каждого $i$-го тела известны:

\begin{itemize}
    \item $m_i$ -- масса тела: $1$
    \item $I_i$ -- момент инерции: $1/6$
    \item $F_{ext,i}$ -- внешняя прикладываемая сила: $(0, -m_i g)$
    \item $\omega_{i}$ -- угловая скорость.
\end{itemize}

В $k$-ом контакте касаются друг друга тела $i$ и $j$. Для каждого контакта известны:

\begin{itemize}
    \item $N_k$ -- сила, с которой тело $i$ давит на тело $j$ по нормали к плоскости касания. 
    \item $\mu$ -- коэффициент трения.
\end{itemize}

В плоскости контакта заведём базис. В двухмерном случае он одномерный, просто нормаль к нормали плоскости касания.....

В этом базисе для каждого контакта заведём:

\begin{itemize}
    \item $v_k$ -- относительная скорость точки контакта как точки первого тела и как точки второго.
    \item $a_k$ -- относительное ускорение точки контакта как точки первого тела и как точки второго.
    \item $f_k$ -- прикладываемая сила трения. Её надо найти.
\end{itemize}

В двумерном случае это всё будут скаляры.

Если $v_k > 0$, тогда однозначно возникает сила трения скольжения:

\[ f_k = -\mu N \cdot \frac{v_k}{|v_k|} \]

Если $v_k = 0$, тогда есть 2 варианта: трение покоя или трение скольжения.
Поэтому силы нужно приложить такие, чтобы выполнялись следующие условия:

\begin{itemize}
    \item $|f_k| \le \mu N$ -- сила трения не превышает своего предела.
    \item $(\mu N - |f_k|) \cdot a_k = 0$ -- относительное ускорение равно нулю
        или сила трения равно своему пределу. 
\end{itemize}

Уравнения движения для тел:

\[
    m_i a_i = F_{ext,i} + \sum_{k} (N_k n_{k,i} + f_k t_k) \\
\]

\[
    I_i \alpha_i = \sum_{k} r_{i,k} \times (N_k n_{k,i} + f_k t_k)
\]

Проще все известные силы (трение скольжения и реакция на контакт) для всех тел представить в виде вектора $F$, а неизвестные силы трения
(покоя или уже перешедшие в скольжение) для всех тел обозначить как вектор $f$.
Матрицу, описывающую "вклад" этих силы к ускорению тела обозначить как $D$. Матрицу масс как $M$, ускорения как $a$ 

Получается:

\[ Ma = F + D^T f \]

\[ a = M^{-1} F + M^{-1} D^T f \]
\[ a = a_{known} + a_{unknown} \]

Чтобы найти относительное ускорение делаем:

\[ a_t = D a \]

\[ a_t = D M^{-1} F + D M^{-1} D^T f \]

Абсолютную величину этих относительных ускорений надо минимизировать:

\[ \min_{f} a_t^T a_t = (D M^{-1} F + D M^{-1} D^T f)^T (D M^{-1} F + D M^{-1} D^T f) \]

\end{document}
